% Generated draft: requires mathematical and source review.
\documentclass[11pt]{article}
\usepackage[margin=25mm]{geometry}
\usepackage{fontspec}
\setmainfont{Latin Modern Roman}
\usepackage{amsmath,amssymb,mathtools,unicode-math}
\setmathfont{Latin Modern Math}
\usepackage{enumitem,needspace,xurl,hyperref}
\setlength{\parindent}{0pt}
\setlength{\parskip}{5pt}
\setlength{\emergencystretch}{3em}
\newcommand{\N}{\mathbb{N}}
\newcommand{\Z}{\mathbb{Z}}
\newcommand{\Q}{\mathbb{Q}}
\newcommand{\R}{\mathbb{R}}
\newcommand{\C}{\mathbb{C}}
\newcommand{\F}{\mathbb{F}}
\newcommand{\A}{\mathbb{A}}
\newcommand{\PP}{\mathbb{P}}
\newcommand{\E}{\mathbb{E}}
\newcommand{\eps}{\varepsilon}
\newcommand{\dd}{\,\mathrm{d}}
\begin{document}
\section*{UnsolvedMath ID: 6200094}
\textbf{Source problem code:} catalogue entry
\subsection{Definitions and mathematical statement}
Let $\mathbb H^n_{\C}$ be complex hyperbolic $n$-space. A standard holomorphic chain in its ideal boundary is the ideal boundary of a complex geodesic, that is, of a totally geodesic holomorphic copy of $\mathbb H^1_{\C}$. Thus chains are distinguished circles in $\partial_\infty\mathbb H^n_{\C}$, with incidence inherited from their intersections and boundary points. The source does not specify which chain pattern or additional structure is to be generalized.

Let $\Gamma<PU(n,1)$ be a lattice, meaning a discrete subgroup of finite covolume. The problem asks for generalizations of these holomorphic chain patterns that can prove rigidity results for embeddings
\[
 \rho:\Gamma\hookrightarrow H
\]
into other higher-rank Lie groups $H$.

This is a research program, not a fully quantified rigidity theorem in the supplied record. The target-group class, meaning of higher rank, boundary or homogeneous geometry carrying the generalized chains, admissible representations, and exact rigidity conclusion are not prescribed. In particular the source does not say whether the embedding must have discrete image, whether cocompactness is required of the lattice, or whether the desired conclusion is extension to an ambient group homomorphism, local rigidity, geometric rigidity, or a classification. None of these alternatives is silently selected. The dimension range and exceptional or degenerate representation cases also remain unspecified.
\subsection{Short English statement}
Develop analogues of the chain geometry on the complex-hyperbolic boundary to obtain rigidity for embeddings of complex-hyperbolic lattices into higher-rank Lie groups. A precise target requires specifying the new boundary geometry, representation hypotheses, and rigidity conclusion.

The supplied summary connects the program to Prasad, Margulis, and Gromov's Kähler/chain rigidity, but gives no exact theorem or reference and says the chain-pattern approach is not developed for all embeddings. Its OPEN-TRIAGE label and the historical 2005 problem are retained with current status explicitly marked for review.
\subsection{Sources}
\begin{itemize}
\item Local UnsolvedMath registry, permanent ID 6200094, source code catalogue entry, ``Boundaries of Groups and Kleinian Groups --- Problem 94,'' proposed by Anna Wienhard, 2005. Dataset curation and metadata: Ulam AI and UnsolvedMath Contributors, CC BY 4.0.
\item Kapovich, \emph{Problems on Boundaries of Groups and Kleinian Groups} (2005), Problem 94, PDF page 23, \url{https://www.math.ucdavis.edu/~kapovich/EPR/problems.pdf}; normalized pdf-text extraction, accessed 2026-07-29.
\item The supplied summary mentions Prasad, Margulis, and Gromov without complete bibliographic references or exact chain-rigidity results. Source-specific provenance marks current status and statement-reuse terms NEEDS\_REVIEW for the public author-hosted PDF, separately from the dataset licence.
\end{itemize}
\end{document}
